\section{Asymptotic behavior of special functions}

\subsection{Leibniz's rule}

We frequently use the following formula, which directly follows from Leibniz's rule.

\begin{Lemma}\label{Leibnizrule}
Let $r$ be a smooth function and $N\in\mathbb{N}\setminus\{0\}$. Then there exist smooth functions~$b_{jN}$ such that
\begin{gather*}
\bigl(\pa_x\circ r(x)^{-1}\bigr)^N=\frac{1}{r(x)^N}\!\sum_{j=0}^Nb_{jN}(x)\pa_{x}^j\qquad \text{with}\quad |b_{jN}(x)|\lesssim \sum_{k=1}^{N-j}\sum_{\substack{\ell_1+\cdots+\ell_k=N-j,\\ \ell_1,\hdots,\ell_k\geq 1}}\prod_{i=1}^k\!\frac{\big|r^{(\ell_i)}(x)\big|}{|r(x)|}.
\end{gather*}
\end{Lemma}

\begin{Remark}
More precisely, we have
\begin{align*}
\bigl(\pa_x\circ r(x)^{-1}\bigr)^N=\frac{1}{r(x)^N}\pa_x^N+\frac{1}{r(x)^N}\sum_{j=0}^{N-1}\Bigg(\sum_{k=1}^{N-j}\sum_{\substack{\ell_1+\cdots+\ell_k=N-j,\\ \ell_1,\hdots,\ell_k\geq 1}}C_{\ell_1\ell_2\hdots \ell_k}^{jN}\prod_{i=1}^{k}\frac{r^{(\ell_i)}(x)}{r(x)} \Bigg)\pa_{x}^j
\end{align*}
with constants \smash{$C_{\ell_1\ell_2\hdots \ell_m}^{jN}>0$}.
\end{Remark}

\subsection{Non-stationary phase theorem with a singular amplitude}

The following lemma is more or less well known and an easy consequence of integration by parts. We give a proof for completeness of the paper.

\begin{Lemma}\label{singularasym}
Let $\m>-1$, $\c\in C^{\infty}((0,\infty))$ which is supported in $x\leq 10$ such that for $\a\in\mathbb{N}$, we have $|\pa_x^{\a}\c(x)|\lesssim |x|^{\m-\a}$ for $x\in (0,\infty)$. Let $f\in C^{\infty}(\R)$ satisfy $\im f(x)\geq 0$, $f'(x)\neq 0$ for~${x\in \supp \c}$ and $f(0)\in \R$. We define
\begin{align*}
b_{\m}(\l):={\rm e}^{-{\rm i}\l f(0)}\int_{0}^{\infty}\c(x){\rm e}^{{\rm i}\l f(x)}{\rm d}x\qquad \text{for}\quad \l\gtrsim 1.
\end{align*}
Then for each $\a\in\mathbb{N}$
\begin{align}\label{singbmues}
|\partial_{\l}^{\a}b_{\m}(\l)|\lesssim \l^{-\m-1-\a},\qquad \l\gtrsim 1.
\end{align}
In particular, $\big|\int_{0}^{\infty}\c(x){\rm e}^{{\rm i}\l f(x)}{\rm d}x\big|\lesssim \l^{-\m-1}$ for $\l\gtrsim 1$.
\end{Lemma}

\begin{proof}
We only need to prove these estimates for sufficiently large $\l$.

First, we prove \eqref{singbmues} for $\a=0$. Let $\chi\in C_c^{\infty}(\R)$ such that $\chi(x)=1$ for $|x|\leq 1/2$ and~${\chi(x)=0}$ for $|x|\geq 1$. Define $\chi_{\l}(x)=\chi(\l x)$ and $\overline{\chi_\l}(x)=1-\chi_{\l}(x)$. Then we have
\begin{align}\label{singbmues1}
\left|\int_{0}^{\infty}\c(x)\chi_{\l}(x){\rm e}^{{\rm i}\l f(x)}{\rm d}x\right|\lesssim \int_{0}^{1/\l}x^{\m}{\rm d}x\lesssim {\l}^{-\m-1}.
\end{align}
On the other hand, the integration by parts yields
\begin{align*}
\left|\int_{0}^{\infty}\c(x)\overline{\chi_{\l}}(x){\rm e}^{{\rm i}\l f(x)}{\rm d}x\right|=\l^{-N}\left|\int_{0}^{\infty}L^N(\c(x)\overline{\chi_{\l}}(x)){\rm e}^{{\rm i}\l f(x)}{\rm d}x\right|,
\end{align*}
where we define $L=\pa_x\circ ({\rm i}f'(x))^{-1}$. By using Leibniz's rule, Lemma \ref{Leibnizrule} and the assumption $f'\neq 0$ on $\supp \c$, we obtain $|L^{N}(\c(x)\overline{\chi_\l}(x))|\lesssim |x|^{\m-N}$. Hence, for $N>\m+1$,
\begin{align}\label{singbmues2}
\left|\int_{0}^{\infty}\c(x)\overline{\chi_{\l}}(x){\rm e}^{{\rm i}\l f(x)}{\rm d}x\right|\lesssim \l^{-N}\int_{\frac{1}{2\l}}^{10}x^{\m-N}{\rm d}x \lesssim \l^{-\m-1}.
\end{align}
The inequalities \eqref{singbmues1} and \eqref{singbmues2} imply \eqref{singbmues} for $\a=0$.

To deal with the case $\a\geq 1$, we write
\begin{align*}
\pa_{\l}^{\a}b_{\m}(\l)={\rm i}^{\a}{\rm e}^{-{\rm i}\l f(0)}\int_0^{\infty}(f(x)-f(0))^{\a}\c(x){\rm e}^{{\rm i}\l f(x)}{\rm d}x.
\end{align*}
By Taylor's theorem, we have $\pa_x^{N}(\c(x)(f(x)-f(0))^{\a})=O\bigl(|x|^{\m+\a-N}\bigr)$ for $x\in \supp \c$. Moreover, $\big|{\rm e}^{-{\rm i}\l f(0)}\big|=1$ due to the assumption $f(0)\in \R$.
Hence the same argument as the case $\a=0$ gives \eqref{singbmues} for $\a\geq 1$.
\end{proof}
